\chapter{نظريات حديثة بديلة للتعلّم القائم على الدوبامين}
\chaptermark{نظريات أخرى في \nt{DOPA}}
\label{sec:dofaalt}
\begin{chaptergoals}
\item كيف تسجل عصبونات \nt{DOPA} ذات الاستجابات اللامتناظرة للأخطاء توزيعَ المكافأة كله؛
\item كيف تنقل بعض عصبونات \nt{DOPA} الأهمية أو الخطر بدل علامة الخطأ؛
\item كيف يحمل سلك واحد خطأً سريعًا يعلّم ومستوى بطيئًا يحدد سرعة الأفعال؛
\item لماذا إشارة \nt{DOPA} حزمة لا عدد واحد، وأي نظرية تنافس الخطأ نفسه.
\end{chaptergoals}

\section{ما الذي يسري في السلك حقًا}

في الفصل~\ref{sec:dofa} كان نظام \nt{DOPA} عندنا سلكًا واحدًا يسري فيه عدد واحد، هو خطأ التنبؤ بالمكافأة $\delta$~\eqref{eq:ctd}. هذه الصورة تفسر الكثير، لكن كلما قيس الدوبامين بدقة أكبر وُجد ما لا ينسجم معها. بعض العصبونات يستجيب للمزعج كما يستجيب للمكافأة؛ وتركيز الدوبامين يرتفع ما دام الحيوان يركض نحو الهدف، مع أنه لا يحدث شيء غير متوقع؛ وعصبونات \nt{DOPA} المختلفة تسلك سلوكًا مختلفًا.

بالنسبة إلى المهندس، كل هذه المكتشفات سؤال واحد: \emph{ما صيغة الإشارة على ناقل البثّ؟} عدد واحد أم متجه؟ إشارة واحدة في السلك أم عدة إشارات موزعة على الترددات؟ هل تتحدث عن المستقبل، مثل $\delta$، أم عن الماضي؟ فيما يلي ستة أجوبة حديثة؛ ولكل منها سنفصل المقيس عن المفترض.

\section{لا عدد واحد بل توزيع}

الخطأ~\eqref{eq:ctd} يعلّم الناقد \emph{متوسط} توقع المكافأة. لكن نصف قطرة عصير مضمونة وقطرة كاملة باحتمال النصف لهما المتوسط نفسه. ولكي نميّز بينهما نحتاج إلى توزيع المكافأة كله.

لنأخذ أبسط مسألة: بعد الإشارة المنبئة تأتي مكافأة حجمها عشوائي $r$. لتكن عصبونات \nt{DOPA} كثيرة، ولكلٍّ منها، وليكن العصبون $i$، تقديره الخاص $V_i$، بحيث يكون خطؤه لحظة المكافأة $\delta_i=r-V_i$. وليأخذ كلٌّ منها الخطأ الموجب بوزن $\alpha_i^+$، والسالب بوزن $\alpha_i^-$. بعد كل مكافأة يتحرك التقدير بمقدار
\begin{equation}
\Delta V_i \;=\; \eta\,\bigl(\alpha_i^+\,[\delta_i]_+ \;-\; \alpha_i^-\,[-\delta_i]_+\bigr).
\label{eq:asymmetric}
\end{equation}
يكفّ التقدير عن التغير حين تتوازن التصحيحات الموجبة والسالبة في المتوسط:
\begin{empheq}[box=\keybox]{equation}
\alpha_i^+\,\bigl\langle[r-V_i]_+\bigr\rangle \;=\; \alpha_i^-\,\bigl\langle[V_i-r]_+\bigr\rangle ,
\qquad
\kappa_i \;=\; \frac{\alpha_i^+}{\alpha_i^++\alpha_i^-} .
\label{eq:expectile}
\end{empheq}
عند $\kappa_i=1/2$ هذا هو المتوسط العادي. وعند $\kappa_i>1/2$ يكون العصبون ”متفائلًا“: تقديره مزاح نحو الطرف الأعلى من التوزيع، وعند $\kappa_i<1/2$ يكون ”متشائمًا“.

مثال: يأتي العصير باحتمال $1/2$، وحجم القطرة $1$. عندئذ تعطي~\eqref{eq:expectile} $\kappa_i\cdot\tfrac12(1-V_i)=(1-\kappa_i)\cdot\tfrac12V_i$، أي $V_i=\kappa_i$ (الشكل~\ref{fig:expectiles}). مجموعة عصبونات بقيم $\kappa_i$ مختلفة تسجّل توزيع المكافأة كما تسجّل مجموعة من المئينات توزيعًا في الإحصاء.

\begin{figure}[t]
\centering
\begin{tikzpicture}[>=Stealth,x=7cm,y=3cm]
\draw[->,gray!70!black] (-0.08,0) -- (1.15,0) node[right,black] {\small المكافأة};
\draw[->,gray!70!black] (-0.08,0) -- (-0.08,0.75) node[left,black] {\small الاحتمال};
\fill[blue!25] (-0.03,0) rectangle (0.03,0.5);
\fill[blue!25] (0.97,0) rectangle (1.03,0.5);
\node[below,font=\small] at (0,0) {$0$};
\node[below,font=\small] at (1,0) {$1$};
\node[left,font=\scriptsize] at (-0.08,0.5) {$1/2$};
\foreach \k/\c/\lab/\h in {0.2/blue!60!black/{متشائم، $\kappa=0.2$}/0.62, 0.5/gray!50!black/{المتوسط، $\kappa=0.5$}/0.42, 0.8/red!70!black/{متفائل، $\kappa=0.8$}/0.22}{
  \draw[very thick,\c] (\k,0) -- (\k,\h);
  \node[above,font=\scriptsize,\c] at (\k,\h) {\lab};}
\end{tikzpicture}
\caption{توزيع المكافأة مسجَّلًا بمجموعة من عصبونات \nt{DOPA}. المكافأة $0$ أو $1$ باحتمال $1/2$ (العمودان الأزرقان)؛ والعصبون ذو النسبة $\kappa$ يصل إلى التقدير $V=\kappa$~\eqref{eq:expectile}.}
\label{fig:expectiles}
\end{figure}

المقيس: لعصبونات \nt{DOPA} المختلفة ”نقاط انقلاب“ مختلفة، أي حجم المكافأة الذي تتحول عنده الاستجابة من ذروة إلى هبوط، والعصبونات ذات اللاتناظر الأكبر بين استجابتيها للخطأ الموجب والسالب تنقلب عند مكافآت أكبر، كما تقتضي~\eqref{eq:expectile}~\cite{Dabney2020}. ومن استجابات مجموعة من العصبونات يمكن استعادة شكل التوزيع. أما أن هذا هو الوظيفة الرئيسية للتشتت، لا أثرًا جانبيًا، ففرضية.

\begin{keyblock}
\begin{proposition}[التوزيع من اللاتناظر]
\label{prop:expectile}
إذا اختلفت عصبونات \nt{DOPA} في النسبة $\kappa_i$ التي تأخذ بها الأخطاء الموجبة، تقاربت تقديراتها إلى نقاط مختلفة~\eqref{eq:expectile} من توزيع المكافأة. فتنقل مجموعة العصبونات لا المتوسط، بل التوزيع، ومعه المخاطرة.
\end{proposition}
\end{keyblock}

\section{لا الخطأ وحده بل الأهمية أيضًا}

وفق صيغة الخطأ~\eqref{eq:ctd} يجب أن يسبب الحدث المزعج، كصدمة أو طعم مرّ، هبوطًا: جاءت النتيجة أسوأ من المتوقع. جزء من عصبونات \nt{DOPA} يفعل ذلك فعلًا، وجزء آخر يستجيب للمزعج \emph{بذروة} كما يستجيب للمكافأة~\cite{Matsumoto2009}. هذه العصبونات لا تخبر بإشارة الخطأ، بل بأن شيئًا \emph{مهمًا} حدث: غير متوقع، قويًا، يستدعي الانتباه~\cite{BrombergMartin2010}.

من المناسب للمهندس أن يفكر في هذا كإشارتين. الأولى الخطأ ذو الإشارة $\delta$: في أي اتجاه تُغيَّر الأوزان. والثانية قيمته المطلقة $|\delta|$ أو جِدّة الحدث: \emph{بأي قوة} تُغيَّر، أي معدل التعلّم؛ فبعد حدث غير متوقع ينبغي التعلّم أسرع. وهي في معناها قريبة من معامل الكسب النورأدرينالي من الفصل~\ref{sec:neurontypes}.

وثمة خيار ثالث: سلك مستقل لمحور مستقل. عصبونات \nt{DOPA} الذاهبة إلى إحدى مناطق اختيار الأفعال لا تستجيب للمكافأة، بل لجِدّة المنبّه وشدته، وهي ضرورية لكي يتعلم الحيوان تجنّب الخطر~\cite{Menegas2018}: ما يسري في السلك ليس ”أفضل أم أسوأ“، بل ”خطر أم لا“.

المقيس: مجموعات مختلفة من عصبونات \nt{DOPA} تستجيب للمزعج وللجديد استجابات مختلفة، والمخارج المختلفة تعلّم أشياء مختلفة. أما كيف يستعمل الدماغ إشارة الأهمية بالضبط، معدلًا للتعلّم، أم إشارة انتباه، أم خطأً مستقلًا على محور الخطر، فمسألة قيد النقاش.

\section{لا التعليم وحده بل الحثّ أيضًا}

للدوبامين أثر فوري أيضًا: كلما زاد، أقبل الحيوان على الفعل برغبة وسرعة أكبر. كيف يتفق هذا مع التعلّم؟

الجواب الهندسي هو \emph{الفصل بالتردد}. الذرى والهبوطات السريعة التي تدوم مئات الميلي ثواني تحمل $\delta$ وتعلّم. والمستوى البطيء الذي يتغير خلال دقائق يحمل مقدارًا آخر، هو متوسط المكافأة في وحدة الزمن $\bar R$~\cite{Niv2007}. ليكن الفعل المنجز في زمن $\tau_a$ يتطلب جهدًا $C/\tau_a$: كلما كان أسرع كان أغلى. لكن كل ثانية تأخير تكلّف $\bar R$، أي قدر المكافأة التي كان يمكن نيلها فيها. الكلفة الكلية $C/\tau_a+\bar R\,\tau_a$ أصغر ما يمكن عند
\begin{empheq}[box=\keybox]{equation}
\tau_a^{*} \;=\; \sqrt{\frac{C}{\bar R}} .
\label{eq:vigor}
\end{empheq}
كلما كانت البيئة أغنى كان الوقت أغلى، وصار الأجدى أن يُعمل بسرعة أكبر: إذا تضاعف $\bar R$ قصر زمن الفعل بـ$\sqrt2\approx1.4$ مرة. ولاحظ أن $\bar R$ هو المكافأة المتوقعة نفسها التي طرحناها في الفصل~\ref{sec:dofa}.

قد يتغير تحرير الدوبامين في موضع الهدف دون أن يتغير تردد النبضات: فالإشارات المحلية عند المخارج نفسها تضبطه~\cite{Mohebi2019}. لكن المكونات البطيئة موجودة في النبضات نفسها أيضًا: في تجارب القسم التالي يُرى الارتفاع التدريجي عند الاقتراب من المكافأة في تردد عصبونات \nt{DOPA} كذلك~\cite{Kim2020}. إذن يتألف المستوى البطيء من مصدرين، تردد النبضات والضبط المحلي للتحرير، ويبدو أن أيهما الغالب يعتمد على المخرج وعلى المهمة.

\begin{keyblock}
\begin{proposition}[إشارتان على سلك واحد]
\label{prop:multiplex}
ينقل نظام \nt{DOPA} مقدارين على الأقل مفصولين في الزمن: خطأً سريعًا $\delta$ يعلّم، ومستوى بطيئًا يحدد ثمن الوقت~\eqref{eq:vigor} ومعه سرعة الأفعال. المقيس أن المكوّنين السريع والبطيء يسلكان سلوكين مختلفين؛ أما أن البطيء يساوي $\bar R$ تحديدًا ففرضية.
\end{proposition}
\end{keyblock}

\section{الارتفاع عند الاقتراب}

إذا ركض الحيوان في ممر نحو مكافأة بعيدة، يرتفع تركيز الدوبامين في منطقة اختيار الأفعال تدريجيًا على مدى ثوانٍ ما دام الهدف يقترب~\cite{Hamid2016}. ووفق الفصل~\ref{sec:dofa} لا ينبغي أن يحدث هذا: كل شيء يسير وفق الخطة، ويجب أن تكون $\delta$ صفرًا.

التفسير الأول: هذا الارتفاع ليس $\delta$، بل التوقع $V$ نفسه، أي إشارة ”الهدف قريب، يستحق الجهد“ من القسم السابق~\cite{Hamid2016}. والتفسير الثاني يحتفظ بـ$\delta$~\cite{Kim2020}. إذا كان التوقع صحيحًا، فإنه عند الأفق $\tau_\gamma$ ينمو في الطريق إلى المكافأة $R$ مثل $V=Re^{-(T-t)/\tau_\gamma}$، ووفق صيغة الخطأ~\eqref{eq:ctd} $\dd V/\dd t-V/\tau_\gamma=0$. لكن لنفترض أن التوقع من بعيد أقل من اللازم، وأن التقدير يتدقق مع الاقتراب. عندئذ ينمو التوقع بانحدار أشد، مثلًا $V=Re^{-(T-t)/\tau_v}$ مع $\tau_v<\tau_\gamma$، و
\begin{empheq}[box=\keybox]{equation}
\delta(t) \;=\; \frac{\dd V}{\dd t}-\frac{V}{\tau_\gamma} \;=\; V\Bigl(\frac{1}{\tau_v}-\frac{1}{\tau_\gamma}\Bigr) \;>\; 0 .
\label{eq:ramp}
\end{empheq}
الخطأ موجب ويزداد مع $V$، وهذا هو الارتفاع التدريجي نفسه (الشكل~\ref{fig:ramp}). عند $\tau_\gamma=4\,\text{s}$ و$\tau_v=1\,\text{s}$ نحصل على $\delta=0.75\,V$ في الثانية. اختُبرت هذه الصيغة في ممر افتراضي بنقل الحيوان لحظيًا إلى نقطة أقرب من الهدف: استجاب الدوبامين بذروة، كما يليق بالمشتقة، لا بمستوى متدرج~\cite{Kim2020}.

\begin{figure}[t]
\centering
\begin{tikzpicture}[>=Stealth,x=1.6cm,y=2.6cm]
\draw[->,gray!70!black] (-4.2,0) -- (0.5,0) node[below left,black] {\small $t$};
\draw[->,gray!70!black] (-4.2,0) -- (-4.2,1.2);
\foreach \x/\l in {-4/{$T-4\,\text{s}$},-2/{$T-2\,\text{s}$},0/{$T\,\text{s}$}}{\draw[gray!60!black] (\x,-0.03) -- (\x,0.03); \node[below,font=\scriptsize] at (\x,0) {\l};}
\draw[thick,densely dashed,gray!60!black] plot[domain=-4:0,samples=40] (\x,{exp(\x/4)});
\draw[very thick,blue!60!black] plot[domain=-4:0,samples=60] (\x,{exp(\x)});
\draw[very thick,red!70!black] plot[domain=-4:0,samples=60] (\x,{0.75*exp(\x)});
\node[anchor=south east,font=\small,gray!50!black] at (-1.3,0.77) {$V$ عند $\tau_v=\tau_\gamma$: $\delta=0$};
\node[right,font=\small,blue!60!black] at (0.05,1.0) {$V$ عند $\tau_v=1\,\text{s}$};
\node[right,font=\small,red!70!black] at (0.05,0.72) {$\delta$ وفق~\eqref{eq:ramp}};
\end{tikzpicture}
\caption{ارتفاع الدوبامين عند الاقتراب من المكافأة بوصفه خطأ تنبؤ، $\tau_\gamma=4\,\text{s}$. الخط المتقطع توقع ينمو تمامًا كما يقتضي الأفق ($\delta=0$)؛ والأزرق توقع ينمو بانحدار أشد؛ والأحمر الخطأ~\eqref{eq:ramp}.}
\label{fig:ramp}
\end{figure}

المقيس: الارتفاع التدريجي موجود، في تركيز الدوبامين وفي تردد نبضات عصبونات \nt{DOPA} معًا~\cite{Kim2020}، والقفزات اللحظية في الموقف تسبب قفزات في الدوبامين. أما أي التفسيرين صحيح، أو هل كلاهما صحيح على مخارج مختلفة، فمسألة قيد النقاش.

\section{النظر إلى الوراء}

كل النظريات السابقة تنظر إلى \emph{الأمام}. وإحدى النظريات الحديثة تعكس الاتجاه~\cite{Jeong2022}. لنفترض أن الدماغ لا يتعلم التنبؤ بالمكافأة من الإشارة المنبئة، بل \emph{يلتفت إلى الوراء} بعد أن ينال المكافأة: أي حدث سبقها أكثر من غيره؟ ومقياس هذه الصلة فرق بين احتمالين:
\begin{empheq}[box=\keybox]{equation}
M \;=\; P(\text{إشارة منبئة قبل المكافأة}) \;-\; P(\text{إشارة منبئة في نافذة عشوائية}) .
\label{eq:retro}
\end{empheq}
إذا ظهرت الإشارة المنبئة كثيرًا دون مكافأة كان الحد الثاني كبيرًا، والصلة ضعيفة. في هذه النظرية لا يُخبر الدوبامين بخطأ التنبؤ، بل بمدى كون الحدث المعيّن \emph{سببًا} محتملًا للمكافأة.

آلية الالتفات إلى الوراء تتطلب ما تتطلبه قاعدة العوامل الثلاثة في الفصل~\ref{sec:dofa}: يجب أن يترك كل حدث أثرًا يذوب، لكي يمكن لحظة المكافأة قراءة ما سبقها. والفرق ليس في المشبك، بل فيما ينقله \nt{DOPA}.

في التجارب التي تتنبأ فيها النظرية~\eqref{eq:retro} والخطأ~\eqref{eq:ctd} بنتائج مختلفة، تبع تحرير الدوبامين في عدد من الحالات الأولى~\cite{Jeong2022}. لكن في تجارب أخرى انزاحت استجابة الدوبامين أثناء التعلّم تدريجيًا من لحظة المكافأة إلى لحظة الإشارة المنبئة، كما يقتضي التعلّم بالخطأ~\eqref{eq:ctd}~\cite{Amo2022}. أي النظريتين تصف الدماغ، لا يُعرف بعد.

\section{الخطأ لا يخص المكافأة وحدها}

التوسيع الأخير يتعلق \emph{بموضوع} الخطأ. إذا استُبدل بالعصير المتوقع عصير آخر بالقيمة نفسها لكن بطعم مختلف، لم تتغير القيمة، والخطأ~\eqref{eq:ctd} يساوي صفرًا. ومع ذلك تستجيب عصبونات \nt{DOPA} لهذا الاستبدال~\cite{Takahashi2017}. والذرى المستحثّة اصطناعيًا للدوبامين تستطيع أن تعلّم الحيوان صلة بين منبّهين محايدين، دون أي مكافأة~\cite{Sharpe2017}. يبدو أن \nt{DOPA} يخبر بخطأ التنبؤ لا بالمكافأة وحدها، بل \emph{بما} سيحدث أيضًا.

صوريًا هذا تعميم بسيط لـ~\eqref{eq:ctd}: بدل دفق مكافأة واحد $r$ تؤخذ مجموعة من سمات ما يجري $\varphi_k$، كالطعم والمكان والصوت، ولكل منها توقعها $M_k$ وخطؤها~\cite{Gardner2018}:
\begin{empheq}[box=\keybox]{equation}
\delta_k(t) \;=\; \varphi_k(t) \;+\; \frac{\dd M_k}{\dd t} \;-\; \frac{M_k}{\tau_\gamma} ,
\qquad k=1,\dots,K .
\label{eq:vectortd}
\end{empheq}
المكافأة ليست إلا واحدة من السمات، ومتجه الأخطاء لا يعلّم ما هو جيد فحسب، بل ما يتبع ماذا أيضًا، أي نموذجًا للعالم.

ويتفق مع هذا عدم تجانس عصبونات \nt{DOPA}: في المهمة الواحدة تحمل عصبونات مختلفة مقادير مختلفة، بعضها مرتبط بالمكافأة، وبعضها بالحركة، وبعضها بموضع الانتباه~\cite{Engelhard2019}.

\begin{keyblock}
\begin{proposition}[حزمة أسلاك بدل سلك]
\label{prop:bundle}
إشارة نظام \nt{DOPA} ليست عددًا واحدًا. إنها موزعة على العصبونات (التوزيع~\eqref{eq:expectile}، والسمات المختلفة~\eqref{eq:vectortd}، ومحور الخطر المستقل) وعلى الزمن (الخطأ السريع والمستوى البطيء~\eqref{eq:vigor}). والخطأ العددي~\eqref{eq:ctd} تقريب أول لهذه الإشارة، كما أن الجامع ذا العتبة تقريب أول للعصبون.
\end{proposition}
\end{keyblock}

\section{الخلاصة}

\begin{table}[t]
\centering
\footnotesize
\setlength{\tabcolsep}{4pt}
\renewcommand{\arraystretch}{1.15}
\begin{tabular}{>{\raggedright}p{2.5cm}>{\raggedright}p{3.2cm}>{\raggedright}p{3.6cm}>{\raggedright\arraybackslash}p{4.2cm}}
\toprule
النظرية & ما يسري في السلك & القياس الرئيسي & ما هو معروف \\
\midrule
خطأ التنبؤ (الفصل~\ref{sec:dofa}) & العدد $\delta$~\eqref{eq:ctd} & ذروة، وانتقال إلى الإشارة المنبئة، وهبوط & مقيس؛ تقريب أول \\
التوزيع & مجموعة $\delta_i$ بلاتناظر مختلف & نقاط انقلاب مختلفة لعصبونات مختلفة & يتفق مع التجربة؛ ودور التشتت فرضية \\
الأهمية & $|\delta|$، الجِدّة؛ محور الخطر & ذرى للمزعج عند جزء من العصبونات & مقيس؛ وكيف يُستعمل قيد النقاش \\
ثمن الوقت & المستوى البطيء $\bar R$ & الدوبامين يسرّع الأفعال؛ والمكوّن البطيء في التحرير عند المخارج وفي تردد النبضات معًا & الفصل بين السريع والبطيء مقيس؛ و$\bar R$ فرضية \\
الارتفاع عند الاقتراب & $V$ أو $\delta$ مع تدقيق التقدير~\eqref{eq:ramp} & ارتفاع تدريجي، وقفزات عند النقل في الممر & الارتفاع مقيس؛ والتفسير قيد النقاش \\
النظر إلى الوراء & مقياس الصلة السببية~\eqref{eq:retro} & تجارب تختلف فيها تنبؤات النظريتين & النتائج متناقضة \\
متجه الأخطاء & $\delta_k$ بحسب السمات~\eqref{eq:vectortd} & استجابة لتغيّر الطعم؛ تعلّم بلا مكافأة & مقيس؛ والصيغة المتجهية فرضية \\
\bottomrule
\end{tabular}
\caption{ما ينقله نظام \nt{DOPA}: الصورة المعيارية في الفصل~\ref{sec:dofa} وتوسيعاتها الحديثة.}
\label{tab:dofaalt}
\end{table}

نظريات الجدول~\ref{tab:dofaalt} لا ينفي بعضها بعضًا: كلها تقريبًا تحتفظ بخطأ التنبؤ أساسًا وتوسّع صيغته. ولا ينافسه جديًا إلا النظر إلى الوراء، والتجارب هي التي ستحسم هذا الخلاف. والخلاصة للمهندس بسيطة: ثمن البثّ في القضية~\ref{prop:broadcastcost} ينخفض إذا كان السلك في الحقيقة أسلاكًا كثيرة بمرسَل إليهم مختلفين.

\section{تمارين}

\begin{exercise}[\lvlA]
\label{ex:07:1}
\emph{نقاط التوزيع لقطرة واحدة.} يأتي عصير حجمه $1$ باحتمال $p=0.2$، وإلا فالمكافأة $0$. جد من~\eqref{eq:expectile} قيم $V_i$ عند $\kappa_i=0.2$ و$0.5$ و$0.8$. كم وسيط هذه المكافأة، وبمَ تختلف النقطة~\eqref{eq:expectile} عن المئين؟
\end{exercise}

\begin{exercise}[\lvlB]
\label{ex:07:2}
\emph{التوزيع من عصبونين.} المكافأة $0$ أو $x$، واحتمال $x$ هو $p$. أفاد عصبونان يعملان بالقاعدة~\eqref{eq:asymmetric} بأن $V(1/2)=0.25$ و$V(0.8)=0.5$. استعد $p$ و$x$ والانحراف المعياري للمكافأة. بمَ سيفيد عصبون عند $\kappa=0.2$؟
\end{exercise}

\begin{exercise}[\lvlB]
\label{ex:07:3}
\emph{نمذجة: التوزيع من اللاتناظر.} حاكِ تسعة عصبونات \nt{DOPA} بقيم $\kappa_i=0.1,\dots,0.9$ والقاعدة~\eqref{eq:asymmetric}، $\alpha_i^+=2\kappa_i$، $\alpha_i^-=2(1-\kappa_i)$، ومكافأة موزعة بانتظام على $[0,1]$. كيف يعتمد تشتت $V_i$ على $\eta$، وأي $\eta$ يلزم لتمييز هذا التوزيع من قطرتين متساويتي الاحتمال $0.25$ و$0.75$؟
\end{exercise}

\begin{exercise}[\lvlA]
\label{ex:07:4}
\emph{ثمن الوقت.} (أ)~اشتق~\eqref{eq:vigor} وجد الكلفة الكلية عند الأمثل. (ب)~قدّر الدماغ $\bar R$ بخطأ قدره $k$ مرة. جد الزيادة النسبية في الكلفة عند $k=2$ و$k=4$. ما الدقة التي يجب أن يخبر بها المستوى البطيء للدوبامين عن $\bar R$؟
\end{exercise}

\begin{exercise}[\lvlB]
\label{ex:07:5}
\emph{الذروة عند النقل.} ينمو التوقع وفق~\eqref{eq:ramp} مع $\tau_v=1\,\text{s}$ و$\tau_\gamma=4\,\text{s}$؛ ويُنقل الحيوان لحظيًا من النقطة $T-2\,\text{s}$ إلى $T-1\,\text{s}$. جد مساحة ذروة $\delta$ كنسبة من $R$، وكم ثانية من الارتفاع التدريجي قبل النقل تحل محلها. ماذا سيُظهر النقل إذا كان الدوبامين يخبر بـ$V$ نفسه؟
\end{exercise}

\begin{exercise}[\lvlB]
\label{ex:07:6}
\emph{تصميم: إشارتان على السلك.} يسري في سلك \nt{DOPA} مستوى ينمو بمعدل $s=1/60$ نبضة في الثانية كل ثانية، وأحداث من $A=3$ نبضات. مشبك أثره $\tau_e=1\,\text{s}$ يطرح من التردد نسخته المنعّمة بـ$\tau_f$. عند أي $\tau_f$ لا يضيع أكثر من $10\,\%$ من الحدث، وتكون مساهمة المستوى خلال زمن الأثر أقل من $10\,\%$ من مساهمة الحدث؟
\end{exercise}

\begin{exercise}[\lvlB]
\label{ex:07:7}
\emph{تصميم تجربة: النظر إلى الوراء ضد الخطأ.} في النافذة تظهر الإشارة المنبئة باحتمال $0.1$، وبعدها المكافأة باحتمال $0.8$. قارن $\Delta P=P(\text{مكافأة}\mid\text{إشارة})-P(\text{مكافأة}\mid\text{لا إشارة})$ و$M$ من~\eqref{eq:retro} إذا (أ)~أُضيفت مكافآت بلا إشارة منبئة، $0.08$ لكل نافذة؛ (ب)~ضوعف تواتر الإشارة المنبئة دون مكافآت.
\end{exercise}

\begin{exercise}[\lvlC]
\label{ex:07:8}
\emph{حزمة الأسلاك وثمن البثّ.} في نموذج التمرين~\ref{ex:06:8} قُسّم $N$ عصبونًا إلى $K$ مجموعة لكل منها سلكها، ويتسرب إلى كل سلك جزء $\rho$ من أخطاء غيره. بيّن أن ثابت التعلّم يساوي $N/K+2+\rho^2N(1-1/K)$. كم محاولة يعطي عند $K\to\infty$ و$N=10^{10}$ و$\rho=0.01$؟
\end{exercise}
